\chapter{データ収集と処理}

\section{専用データ収集ソフトウェア}

本プロジェクトは、数本の動画を暫定保存するだけでなく、ALOHA 専用の収集ソフトウェアを開発した。確認できた機能は、双腕遠隔操作、連続収集、非同期保存、破棄後の再収集、キーボード・足踏み・遠隔トリガ、ランダム開始点、初期姿勢へ戻る/戻らない二つの手順、連続回転関節、軌跡と動画の同時保存、ハードウェア符号化が使えない場合の代替処理、軌跡完全性検査、ロボット健全性検査、安全停止である。\ev{E02}

\plain{収集ソフトウェアはデータ生産ラインに相当する。収集担当者は連続記録し、誤りを見つければ直ちに破棄して再収集できる。動画保存をバックグラウンド化し、各軌跡後の待ち時間を減らす。}

本監査が確認したのは 2026 年 7 月時点の機能である。ソフトウェアは継続的に更新されているため、2026 年 5 月の本番学習データすべてに現在の機能が適用されたとは限らない。

\section{データ編集・分析センター}

編集センターは ALOHA に限らず複数ロボットのデータを集中管理する。読み取り専用監査では、有効 ALOHA プロジェクト 51 件、有効軌跡 2,413 本、2,907,804 フレーム、58,156.08 秒（約 16.15 時間）を確認した。プロジェクト集計と軌跡ごとの再集計は一致した。現行の有効プロジェクト外にある 562 軌跡はこの合計から除外した。\ev{E03}

プラットフォームは、プロジェクト・軌跡閲覧、単一軌跡の多視点同期表示、単体・一括ラベル、データセット生成、サブタスク生成、処理状況追跡、停止、再実行に対応し、複数種のロボットデータを扱える。これにより、各計算機に散在するファイルを、監査・バージョン管理・再学習可能なデータ資産へ変換した。

\begin{table}[H]
\centering
\small
\caption{データ編集・分析センターの主要工程}
\begin{tabularx}{\textwidth}{>{\japanesesans\bfseries}p{27mm}p{47mm}X}
\toprule
工程 & 実装済み機能 & 学習への直接的な価値\\
\midrule
プロジェクト管理 & ロボット、タスク、収集ロット単位で整理 & 出所と版を区別し、異なる集計基準の混在を防ぐ\\
軌跡閲覧 & 単一軌跡と同期多視点映像を表示 & 遮蔽、向き、キャップ状態、双腕協調の問題を発見\\
ラベル編集 & 個別・一括ラベル、サブタスク分割 & キャップ有無、向き、段階、異常を学習条件へ変換\\
データ生成 & 審査済み軌跡から固定版を生成 & 学習入力の追跡性と再現性を確保\\
処理タスク & 生成状態を追跡し、停止・再実行を支援 & 大規模処理の状態管理と失敗復旧を実現\\
複数ロボット対応 & 同一基盤で異なるロボットデータを管理 & ALOHA 以外の後続タスクにも基盤を再利用\\
\bottomrule
\end{tabularx}
\end{table}

\plain{ラベルは画面を見栄えよくするためではない。「軌跡で何が起きたか」を学習条件に変換するために付与する。キャップ有無、口部方向、作業段階が明確でなければ、モデルは条件に応じた行動選択を学びにくい。}

\figfull[.96\textwidth]{data_scope.pdf}{データ資産の三つの集計範囲。基盤全体には約 16.15 時間、実機モデルの重複なし学習部分には約 4.89 時間、フィルタ後の実学習部分には約 4.69 時間がある。三つを混同しない。}

\section{実機モデルが使用したデータ}

\begin{table}[H]
\centering
\small
\caption{実機運用モデルの学習データ}
\begin{tabularx}{\textwidth}{>{\japanesesans\bfseries}p{31mm}p{39mm}X}
\toprule
指標 & 監査値 & 説明\\
\midrule
重複しないデータ保管単位 & 25 & 学習設定の 29 サンプリング項目には重み付け用の重複を含む\\
重複なし軌跡 & 1,051 & 各軌跡は人による双腕示教\\
重複なしフレーム & 879,852 & 50 fps 換算で約 4.89 時間\\
学習対象フレーム & 844,102 & 約 4.69 時間。35,750 フレームは除外\\
重み付き曝露量 & 1,495 軌跡相当 & 重点データを反復抽出。新規収集数ではない\\
原画像 & 4 視点、640$\times$480、50 fps & 実機モデルは上方、左手首、右手首の 3 視点を使用\\
学習画像 & 3 視点、224$\times$224 & 学習時に統一リサイズ\\
状態・行動 & 各 14 次元 & 双腕各 6 関節と 1 グリッパ\\
言語指示 & あり & データ保管単位ごとのタスク記述から取得\\
報酬・成功ラベル & なし & 基礎模倣データから成功率は直接算出不可\\
データ分割 & 学習のみ & 独立した検証・テスト集合なし\\
\bottomrule
\end{tabularx}
\end{table}

公開データは Hugging Face の
\href{https://huggingface.co/lyl472324464}{huggingface.co/lyl472324464}
に置かれている。Dataset Viewer は行単位のプレビュー、フィルタ、統計、画像・動画項目の確認に対応する\cite{hf_dataset_viewer}。本報告書は、確定版データの説明情報、集計表、動画を相互確認し、ページ名だけには依存しない。

\figfull[.85\textwidth]{sampling_exposure.pdf}{重点条件の反復サンプリング。1 サイクル当たり 1,495 軌跡相当を提示し、重複なし 1,051 軌跡を上回る。増えるのは難条件の提示回数であり、収集軌跡数ではない。}

\section{条件別データの範囲}

\figfull[.95\textwidth]{condition_coverage.pdf}{実機学習データに含まれる条件。方向、キャップなし、復帰、水入り、反転、自由回転キャップを対象に追加収集と重み付けを行った。名称キーワードによる分類は重複を許し、成功率を表さない。}

方向変化 404 軌跡、キャップなし 158、復帰 135、水入り 136、反転 127、自由回転キャップ 111 を確認した。実機学習では自由回転キャップデータを 5 回反復抽出しており、難条件への学習曝露を意図的に増やしている。\ev{E04,E11}

\figfull[.92\textwidth]{baseline_episode_length_distribution.pdf}{1,051 軌跡の長さ分布。中央値は 12.2 秒、中央 80\% は約 4.9--28.4 秒であり、異なる長さの作業を含む。可読性のため極端な長尾は 99 パーセンタイルで表示を打ち切った。}

\section{実示教フレーム}

\figfull[\textwidth]{training_demonstration_keyframes.png}{4 条件の実示教フレーム。各条件で中央値に近い軌跡を選び、20\%、50\%、80\% 時点を示す。通常、方向変化、キャップなし、反転を含む。人の示教データであり、自律評価結果ではない。}

\section{品質検査、正規化、拡張}

\begin{table}[H]
\centering
\small
\caption{数値品質の全数検査}
\begin{tabularx}{\textwidth}{>{\japanesesans\bfseries}p{42mm}p{28mm}X}
\toprule
検査項目 & 結果 & 範囲\\
\midrule
状態・行動の非有限値 & 0 & 全数値レコード\\
全要素がゼロの状態レコード & 0 & 全数値レコード\\
全要素がゼロの動作レコード & 0 & 全数値レコード\\
軌跡内時刻の逆転・停止 & 0 軌跡 & 保存時刻に基づく\\
軌跡内フレーム番号欠落 & 0 軌跡 & 保存番号に基づく\\
完全一致数値軌跡群 & 0 群 & 視覚的意味の重複とは異なる\\
破損動画 & 全数未確認 & 代表動画のみ抽出・復号\\
学習・テスト漏洩 & 検査不可 & 学習分割しか存在しない\\
\bottomrule
\end{tabularx}
\end{table}

各状態・行動次元には 1、99 パーセンタイルを保存する。スカラー $x$ の分位正規化は次式で表せる。
\[
\tilde{x}=2\,\frac{\operatorname{clip}(x,q_{01},q_{99})-q_{01}}
{q_{99}-q_{01}}-1.
\]
通常値を $[-1,1]$ に写像し、外れ値の影響を抑える。出力動作は逆変換する。\ev{E09} 画像にはランダム切出し、拡大縮小、回転、色変動を用いるが、全データ保管単位への適用確率は当該学習記録で確認する必要がある。

\limitbox{データは学習分割のみで、独立検証・テスト集合がない。全動画を全フレーム復号しておらず、キャップ有無、ボトル方向、作業段階のフレーム単位ラベルもない。これらは次年度の信頼できる評価を妨げる直接要因である。}

\begin{table}[H]
\centering
\small
\caption{データ監査で確認できる範囲}
\begin{tabularx}{\textwidth}{>{\japanesesans\bfseries}p{32mm}X}
\toprule
区分 & 結論\\
\midrule
確認可能 & 規模、固定版、軌跡長、3 学習カメラ、14 次元状態・行動、数値完全性、重点サンプリング、言語指示は追跡可能。\\
部分確認 & 条件名は収集意図を示すがフレーム単位意味ラベルはない。代表動画は復号できるが全動画ではない。\\
確認不可 & 独立検証・テストがなく漏洩不在を証明できない。成功ラベルがなく自律成功率を算出できない。\\
\bottomrule
\end{tabularx}
\end{table}
